% attinormativi-esempi.tex
%
% Documento di esempio autonomo e compilabile per la classe attinormativi.
% Standalone compilable example for the attinormativi class.
%
% Versione 1.2 — copre tutte le funzionalità della classe, inclusa
% l'opzione accessible (PDF/UA-1, Direttiva UE 2016/2102).
%
% ---------------------------------------------------------------------------
% ISTRUZIONI DI COMPILAZIONE
%
%   MODALITÀ STANDARD (senza accessibilità):
%     pdflatex attinormativi-esempi.tex   (tre passate)
%     xelatex  attinormativi-esempi.tex
%     lualatex attinormativi-esempi.tex
%
%   MODALITÀ ACCESSIBLE (PDF/UA-1):
%     Richiede OBBLIGATORIAMENTE LuaLaTeX e TeX Live 2023+.
%     Decommentare il blocco \DocumentMetadata e l'opzione "accessible"
%     nella \documentclass (Sezione A), quindi compilare con:
%       lualatex attinormativi-esempi.tex
%       lualatex attinormativi-esempi.tex
%       lualatex attinormativi-esempi.tex   (tre passate per tagpdf)
%
%     Verifica del PDF prodotto:
%       PAC 2024 (PDF Accessibility Checker, gratuito, Windows)
%       Adobe Acrobat Pro: Strumenti → Accessibilità → Verifica completa
%
% Il file attinormativi.cls deve trovarsi nella stessa directory oppure
% in una directory riconosciuta da TEXMF.
%
% ---------------------------------------------------------------------------
% VARIANTI DA PROVARE PER LE OPZIONI STANDARD
%
% 1. Font sans-serif:
%      \documentclass[bozza,sansserif,12pt]{attinormativi}
%
% 2. Impaginazione fronte/retro:
%      \documentclass[bozza,12pt,duplex]{attinormativi}
%
% 3. Margini personalizzati:
%      \documentclass[bozza,12pt,left=40mm,right=25mm,top=30mm,bottom=30mm]{attinormativi}
%
% Per impostazione predefinita (questa compilazione base):
%   - font serif (Latin Modern Roman)
%   - margini standard (left=35mm, right=30mm, top=25mm, bottom=25mm)
%   - impaginazione solo fronte (simplex/oneside)
%
% Licenza: LPPL 1.3c
% ---------------------------------------------------------------------------

% ---------------------------------------------------------------------------
% SEZIONE A — Opzioni della classe
%
% Attivare/disattivare le opzioni di seguito per verificare le singole
% funzionalità.  Il set di default mostra: bozza + sansserif + simplex.
%
% Per testare sansserif: aggiungere "sansserif" alle opzioni.
% Per testare duplex:    sostituire "simplex" con "duplex".
% Per testare margini:   aggiungere, ad es., "left=40mm,right=25mm".
%
% ---------------------------------------------------------------------------
% PER TESTARE L'ACCESSIBILITÀ PDF/UA-1 (opzione accessible):
%
%   1. Decommentare il blocco \DocumentMetadata qui sotto.
%   2. Decommentare "accessible" nella \documentclass.
%   3. Compilare SOLO con LuaLaTeX (tre passate).
%   4. Verificare il PDF con PAC 2024 o Adobe Acrobat Pro.
%
%   NOTA: \DocumentMetadata deve precedere \documentclass. Se si usa
%   la modalità standard (senza accessible), lasciarlo commentato.
% ---------------------------------------------------------------------------

%% ── Decommentare il blocco seguente per attivare la modalità accessible ──
%\DocumentMetadata{
%  pdfversion  = 1.7,
%  pdfstandard = UA-1,
%  lang        = it-IT,
%}
%% ────────────────────────────────────────────────────────────────────────

\documentclass[%
  bozza,          %% modalità bozza: attiva note redazionali e banner
  sansserif,      %% font: sansserif (Latin Modern Sans). Sostituire con "serif".
  simplex,        %% impaginazione: solo fronte. Sostituire con "duplex".
                  %%   Alias accettati: oneside / twoside.
  12pt            %% corpo del testo
  %% accessible,  %% <-- Decommentare (con \DocumentMetadata sopra) per PDF/UA-1
]{attinormativi}
%
% ---------------------------------------------------------------------------
% Variante margini — decommentare la riga seguente per testarla:
% \documentclass[bozza,serif,simplex,12pt,left=40mm,right=25mm,top=30mm,bottom=28mm]{attinormativi}
% ---------------------------------------------------------------------------

% ---------------------------------------------------------------------------
% SEZIONE B — Pacchetti opzionali dell'utente
%
% cleveref va caricato nel preambolo DOPO \documentclass{attinormativi}
% e PRIMA di \begin{document}.  La classe rileva automaticamente la sua
% presenza e attiva le definizioni \crefname italiane in \AtBeginDocument.
% ---------------------------------------------------------------------------

\usepackage{cleveref}
\usepackage{hologo}   %% per \hologo{LuaLaTeX} nella Sezione H

% ---------------------------------------------------------------------------
% SEZIONE C — Metadati dell'atto
% ---------------------------------------------------------------------------

\tipoatto{Regolamento}
\numeroatto{3}
\dataatto{1° maggio 2026}
\autoritaatto{Comune di Esempiopoli}
\titoloatto{Regolamento per la promozione delle energie rinnovabili
e delle comunità energetiche}
\oggettoatto{Disciplina delle modalità di costituzione, adesione
e gestione delle comunità energetiche rinnovabili nel territorio comunale,
in attuazione del decreto legislativo 8~dicembre~2021, n.~199.}

% ---------------------------------------------------------------------------
% Per il test dell'alias \maketitle si veda la Sezione D, punto 1.
% ---------------------------------------------------------------------------

\begin{document}



% ===========================================================================
% SEZIONE D — Frontespizio
%
% Punto 1: \maketitle come alias di \frontespizioatto.
%
% I due comandi sono equivalenti. In questa compilazione si usa \maketitle
% per verificare che l'alias funzioni.  Per verificare l'originale basta
% sostituire \maketitle con \frontespizioatto: il risultato deve essere
% identico.
% ===========================================================================

\maketitle            %% TEST alias: equivalente a \frontespizioatto

% ── Banner di bozza (visibile solo con l'opzione "bozza") ──────────────────

\bozzainfo

% ── Indici di apertura ─────────────────────────────────────────────────────

\tableofcontents
\indicearticoli

% ===========================================================================
% PREAMBOLO
% ===========================================================================

\begin{preambolo}
Visto il decreto legislativo 8~dicembre~2021, n.~199, recante attuazione
della direttiva (UE) 2018/2001 sulla promozione dell'uso dell'energia da
fonti rinnovabili;\\
Visto il decreto-legge 1°~marzo~2022, n.~17, convertito con modificazioni
dalla legge 27~aprile~2022, n.~34;\\
Vista la deliberazione del Consiglio comunale n.~45 del 10~marzo~2026,
con la quale si è approvato il Piano d'azione per l'energia sostenibile;\\
Ritenuto necessario disciplinare a livello locale la costituzione e il
funzionamento delle comunità energetiche rinnovabili;\\
Il Consiglio comunale adotta il seguente regolamento.
\end{preambolo}

% ===========================================================================
\titolo{Disposizioni generali}
% ===========================================================================

\articolo{Finalità e ambito di applicazione}[art:finalita]

\comma{Il presente regolamento disciplina la costituzione, l'adesione e
la gestione delle comunità energetiche rinnovabili (di seguito: CER) nel
territorio del Comune di Esempiopoli, in attuazione del decreto legislativo
8~dicembre~2021, n.~199.}[c:finalita:1]

\comma{Le disposizioni del presente regolamento si applicano alle CER
costituite nella forma di ente del terzo settore, di cooperativa, di
associazione o di altro soggetto di diritto privato senza scopo di lucro
prevalente, la cui area di riferimento ricada, anche parzialmente, nel
territorio comunale.}

\comma{Sono esclusi dall'ambito di applicazione del presente regolamento:
\begin{lettere}
  \item i gruppi di autoconsumo collettivo di cui all'articolo~42 del
        decreto legislativo 8~dicembre~2021, n.~199;
  \item le comunità energetiche disciplinate da normativa regionale
        speciale incompatibile con le presenti disposizioni.
\end{lettere}
}

% ---------------------------------------------------------------------------

\articolo{Definizioni}[art:definizioni]

\comma{Ai fini del presente regolamento si intende per:
\begin{definizioni}
  \definizione{Comunità energetica rinnovabile}{Il soggetto giuridico
    autonomo, costituito ai sensi dell'articolo~31 del decreto legislativo
    8~dicembre~2021, n.~199, basato sulla partecipazione aperta e
    volontaria, che produce, consuma, condivide e vende energia elettrica
    da fonti rinnovabili.}[def:cer]

  \definizione{Impianto di produzione}{L'insieme di macchine, apparecchi,
    dispositivi e strutture destinate alla produzione di energia elettrica
    da fonti rinnovabili, allacciati alla rete di distribuzione in bassa
    o media tensione nella medesima zona di mercato.}[def:impianto]

  \definizione{Membro produttore}{Il soggetto partecipante alla CER che
    detiene, a qualsiasi titolo, uno o più \defref{def:impianto} conferiti
    nella comunità.}[def:membro-produttore]

  \definizione{Membro consumatore}{Il soggetto partecipante alla CER che
    preleva energia dalla rete senza disporre di un proprio
    \defref{def:impianto}.}[def:membro-consumatore]

  \definizione{Energia condivisa}{La quantità di energia elettrica,
    espressa in chilowattora, prodotta e contemporaneamente consumata
    dai membri della CER nell'arco della medesima ora, calcolata ai sensi
    della vigente normativa tecnica del Gestore dei Servizi Energetici.}[def:energia-condivisa]

  \definizione{Sportello CER}{Lo sportello unico istituito dal Comune per
    l'assistenza ai soggetti interessati alla costituzione o all'adesione
    a una comunità energetica rinnovabile.}[def:sportello]
\end{definizioni}
}

\notaredaz{Verificare la coerenza delle definizioni con le linee guida
del GSE pubblicate nel febbraio~2024 e con gli aggiornamenti normativi
eventualmente intervenuti prima dell'adozione definitiva.}

% ===========================================================================
% SEZIONE E — Test cleveref
%
% I comandi \cref{} e \Cref{} devono produrre il nome del livello
% strutturale in italiano, senza che l'utente debba scriverlo a mano.
%
% Risultati attesi (cleveref + definizioni italiane della classe):
%   \cref{art:finalita}       => "articolo 1"
%   \Cref{art:finalita}       => "Articolo 1"       (a inizio frase)
%   \cref{c:finalita:1}       => "comma 1"
%   \cref{def:cer}            => "definizione a"
%   \cref{all:statuto}        => "allegato A"
%   \cref{art:finalita,art:requisiti}   => lista automatica ("articoli 1 e 3")
%
% Il blocco seguente è visibile solo in modalità bozza; in produzione
% non occupa spazio ma i label e i riferimenti rimangono attivi.
% ===========================================================================

\notaredaz{%
  \textbf{TEST cleveref — riferimenti automatici.}\\[0.4ex]
  I valori tra parentesi sono i risultati attesi; quelli a sinistra
  del trattino sono prodotti da cleveref.
  \begin{itemize}
    \item \cref{art:finalita}
          \quad\textrightarrow\quad \textit{atteso:} articolo~1
    \item \Cref{art:finalita}
          \quad\textrightarrow\quad \textit{atteso:} Articolo~1
    \item \cref{c:finalita:1}
          \quad\textrightarrow\quad \textit{atteso:} comma~1
    \item \cref{def:cer}
          \quad\textrightarrow\quad \textit{atteso:} definizione~a
    \item \cref{all:statuto}
          \quad\textrightarrow\quad \textit{atteso:} allegato~A
    \item \cref{art:finalita,art:requisiti}
          \quad\textrightarrow\quad \textit{atteso:} articoli~1 e~3
    \item \crefrange{art:finalita}{art:definizioni}
          \quad\textrightarrow\quad \textit{atteso:} articoli~1--2
  \end{itemize}
}

% ===========================================================================
\titolo{Costituzione e riconoscimento delle CER}
% ===========================================================================

\articolo{Requisiti per il riconoscimento comunale}[art:requisiti]

\comma{Il Comune riconosce, ai fini dell'applicazione dei benefici previsti
dal presente regolamento, le CER che soddisfano congiuntamente i seguenti
requisiti:
\begin{lettere}
  \item siano formalmente costituite e registrate secondo le modalità
        previste dalla normativa vigente;
  \item abbiano sede legale o operativa nel territorio comunale;
  \item abbiano adottato uno statuto conforme allo schema di cui
        al \cref{all:statuto};      %% TEST: cleveref su allegato
  \item abbiano presentato la documentazione di cui
        al \cref{art:domanda}.      %% TEST: cleveref su articolo
\end{lettere}
}[c:requisiti:1]

\comma{Il riconoscimento è subordinato alla verifica, da parte
dello \defref{def:sportello}, della disponibilità di almeno un
\defref{def:impianto} allacciato alla rete nella zona di mercato
di riferimento.}

% ---------------------------------------------------------------------------

\articolo{Procedura di riconoscimento}[art:domanda]

\comma{La CER richiedente presenta allo \defref{def:sportello} apposita
domanda, redatta su modello predisposto dall'ufficio, corredata dalla
seguente documentazione:
\begin{numeri}
  \item copia dell'atto costitutivo e dello statuto;
  \item elenco dei membri fondatori con indicazione della qualifica di
        \defref{def:membro-produttore} o \defref{def:membro-consumatore};
  \item documentazione tecnica relativa agli impianti di produzione
        conferiti, con indicazione della potenza nominale installata;
  \item dichiarazione sostitutiva di atto notorio attestante l'insussistenza
        di cause ostative al riconoscimento.
\end{numeri}
}

\comma{Lo sportello conclude il procedimento entro trenta giorni dalla
ricezione della domanda completa. In caso di documentazione incompleta,
il responsabile del procedimento richiede, entro dieci giorni, le
integrazioni necessarie; il termine si sospende fino alla ricezione
delle stesse.}

\comma[bis]{Decorso il termine di cui al comma~2 senza che sia intervenuto
un provvedimento espresso, la domanda si intende accolta ai sensi
dell'articolo~20 della legge 7~agosto~1990, n.~241.}

\notaredazinline{verificare compatibilità con le procedure SUAP}

% ===========================================================================
\titolo{Incentivi e benefici comunali}
% ===========================================================================

\capo{Benefici di carattere urbanistico}

\articolo{Procedura semplificata per gli impianti FER}[art:procedura-semplificata]

\comma{Gli impianti fotovoltaici e le relative opere di connessione,
destinati ad alimentare una CER riconosciuta ai sensi
del \cref{art:requisiti},    %% TEST: cleveref su articolo (testo corrente)
beneficiano della procedura di comunicazione semplificata di cui
all'articolo~6 del decreto legislativo 3~marzo~2011, n.~28, a prescindere
dalla potenza nominale installata, entro i limiti consentiti dalla
normativa nazionale.}

\comma{Il responsabile del procedimento comunica l'esito della verifica
entro quindici giorni dalla presentazione della comunicazione.}

\capo{Benefici di carattere fiscale}

\articolo{Riduzione dell'imposta municipale}[art:imu]

\comma{Gli immobili sui quali sono installati impianti di produzione
conferiti a una CER riconosciuta beneficiano di una riduzione dell'aliquota
dell'imposta municipale propria (IMU) nella misura stabilita annualmente
dal Consiglio comunale con la delibera di approvazione delle aliquote,
nei limiti previsti dalla legislazione vigente.}

\comma{La riduzione di cui al comma~1 cessa automaticamente in caso di
revoca del riconoscimento ai sensi del \cref{art:revoca}.}
%% TEST: cleveref — articolo non ancora dichiarato al momento del \comma;
%%       cleveref risolve il riferimento alla terza compilazione.

% ===========================================================================
\titolo{Revoca del riconoscimento e controlli}
% ===========================================================================

\articolo{Revoca del riconoscimento}[art:revoca]

\comma{Il Comune revoca il riconoscimento qualora la CER:
\begin{lettere}
  \item perda uno o più requisiti di cui al \cref{art:requisiti};
  \item non trasmetta la relazione annuale di cui al \cref{art:relazione}
        entro i termini stabiliti;
  \item violi in modo grave o reiterato le disposizioni del presente
        regolamento.
\end{lettere}
}

\comma{Prima della revoca, il responsabile del procedimento comunica
alla CER i motivi ostativi, assegnando un termine non inferiore a quindici
giorni per la presentazione di memorie e documenti ai sensi dell'articolo~10
della legge 7~agosto~1990, n.~241.}

% ---------------------------------------------------------------------------

\articolo{Obblighi di rendicontazione}[art:relazione]

\comma{Le CER riconosciute trasmettono annualmente allo \defref{def:sportello},
entro il 31~marzo di ciascun anno, una relazione contenente:
\begin{lettere}
  \item i dati di produzione, consumo e \defref{def:energia-condivisa}
        dell'anno precedente, estratti dal sistema di misura del Gestore
        della Rete di Distribuzione;
  \item l'elenco aggiornato dei membri, con separata indicazione dei
        \defref{def:membro-produttore} e dei \defref{def:membro-consumatore};
  \item la quota dei benefici economici derivanti dall'incentivo
        sull'\defref{def:energia-condivisa} destinata a utenze in condizioni
        di disagio economico.
\end{lettere}
}

% ===========================================================================
\titolo{Disposizioni finali}
% ===========================================================================

\normatransitoria{I soggetti già operanti, alla data di entrata in vigore
del presente regolamento, come comunità energetiche rinnovabili nel
territorio comunale, sono tenuti a presentare domanda di riconoscimento
ai sensi del \cref{art:domanda} entro novanta giorni dall'entrata in
vigore. Nelle more del procedimento, continuano ad applicarsi le disposizioni
previgenti.}
%% TEST cleveref in disposizione finale tipizzata.

\abrogazioni{Sono abrogate tutte le disposizioni regolamentari comunali
incompatibili con il presente regolamento. In particolare, è abrogato
l'articolo~12 del regolamento per la disciplina degli impianti tecnologici
approvato con delibera consiliare n.~18/2019.}

\entrataInVigore{Il presente regolamento è pubblicato all'albo pretorio
online del Comune di Esempiopoli ed entra in vigore il quindicesimo giorno
successivo alla data di pubblicazione.}

% ===========================================================================
% ALLEGATI
% ===========================================================================

\allegato{Schema di statuto della comunità energetica rinnovabile}[all:statuto]

Il presente schema costituisce un modello orientativo. Le CER richiedenti
possono adattarlo alle proprie specificità organizzative nel rispetto
dei requisiti inderogabili stabiliti dal \cref{art:requisiti}.
%% TEST cleveref in corpo allegato.

\articolo[1]{Denominazione e sede}
\comma{È costituita la comunità energetica rinnovabile denominata
«\textit{CER [nome]}», con sede legale in Esempiopoli (XX), via
\underline{\hspace{4cm}}, n.~\underline{\hspace{1cm}}.}

\articolo[2]{Oggetto e finalità}
\comma{La comunità ha per oggetto la produzione, il consumo, la condivisione
e la vendita di energia elettrica da fonti rinnovabili, senza scopo di
lucro prevalente, nel rispetto della normativa vigente.}

% ── ── ── ──

\allegato{Modulo di domanda di riconoscimento}[all:modulo]

\noindent\textbf{Al Comune di Esempiopoli}\\
\textbf{Sportello CER}\\[1em]

Il/La sottoscritto/a \underline{\hspace{5cm}}, nato/a a
\underline{\hspace{3.5cm}} il \underline{\hspace{2cm}}, C.F.\
\underline{\hspace{4cm}}, in qualità di legale rappresentante della
CER \underline{\hspace{4cm}}, con sede in \underline{\hspace{4cm}},

\bigskip
\noindent\textbf{CHIEDE}

\noindent il riconoscimento della comunità energetica rinnovabile
ai sensi del \cref{art:requisiti} del Regolamento comunale per
la promozione delle energie rinnovabili e delle comunità energetiche,
allegando la documentazione di cui al \cref{art:domanda}.
%% TEST: cleveref nel modulo allegato (riferimenti incrociati a distanza).

\bigskip
\noindent Esempiopoli, lì \underline{\hspace{3cm}}

\bigskip
\noindent Firma del legale rappresentante: \underline{\hspace{5cm}}

% ===========================================================================
% SEZIONE F — Riepilogo test funzionalità v1.1
%
% Visibile solo in modalità bozza.
% Questo blocco riepiloga i risultati attesi e funge da checklist
% per chi verifica manualmente il PDF prodotto.
% ===========================================================================

\notaredaz{%
  \textbf{CHECKLIST VERIFICA FUNZIONALITÀ v1.2}\\[0.5ex]
  %
  \textbf{1 — Opzione font.}
  Il documento corrente usa \texttt{sansserif} (Latin Modern Sans).
  Ricompilare con l'opzione \texttt{serif} e verificare che:
  (a)~il font del testo cambi in Latin Modern Roman;
  (b)~la giustificazione e la sillabazione rimangano attive;
  (c)~non compaiano overfull box sistematici.\\[0.4ex]
  %
  \textbf{2 — Alias \texttt{\textbackslash maketitle}.}
  Il frontespizio è stato generato con \texttt{\textbackslash maketitle}.
  Verificare che sia identico a quello prodotto da
  \texttt{\textbackslash frontespizioatto} sostituendo il comando
  e ricompilando.\\[0.4ex]
  %
  \textbf{3 — Supporto cleveref.}
  Nel blocco \emph{TEST cleveref} in apertura del Titolo II e negli
  altri punti marcati \texttt{\%\% TEST cleveref} verificare che:
  (a)~\texttt{\textbackslash cref} produca la forma minuscola corretta
  (\textit{articolo}, \textit{comma}, \textit{definizione}, \textit{allegato});
  (b)~\texttt{\textbackslash Cref} produca la forma maiuscola a inizio frase;
  (c)~i rimandi multipli (\texttt{\textbackslash cref\{art:finalita,art:requisiti\}})
  e gli intervalli (\texttt{\textbackslash crefrange}) siano formattati
  automaticamente.\\[0.4ex]
  %
  \textbf{4 — Opzioni margini e fronte/retro.}
  Ricompilare con ciascuna delle seguenti varianti dichiarate in testa
  al file e verificare visivamente il PDF:
  \begin{itemize}
    \item \texttt{duplex}: le pagine pari e dispari devono avere margini
          asimmetrici (inner $>$ outer); il primo capitolo deve iniziare
          su pagina dispari.
    \item \texttt{left=40mm,right=25mm}: la colonna di testo deve spostarsi
          verso destra rispetto al default.
    \item \texttt{top=30mm,bottom=30mm}: il bianco superiore e inferiore
          deve aumentare rispetto al default di 25mm.
  \end{itemize}
  \textbf{5 — Accessibilità PDF/UA-1 (opzione \texttt{accessible}).}
  Per i test specifici dell'accessibilità vedere la \textbf{Sezione~H}
  in fondo al documento. Quella sezione contiene la checklist completa
  con le istruzioni per la verifica con PAC~2024 e la spiegazione
  di ogni elemento taggato.
}

% ===========================================================================
% SEZIONE G — Test numerazione manuale degli articoli
%
% Questa sezione verifica tutti i casi d'uso del parametro opzionale [N]
% su \articolo, \comma e sulle disposizioni finali tipizzate.
%
% Struttura del test:
%   Titolo VI   — Articoli con numerazione automatica (baseline)
%   Titolo VII  — Salto in avanti:    \articolo[20]{...}
%   Titolo VIII — Salto all'indietro: \articolo[5]{...}
%   Titolo IX   — Ripresa automatica dopo interruzione manuale
%   Titolo X    — Numero manuale con label e verifica riferimenti
%   Titolo XI   — Comma manuale e ripresa automatica del contatore
%   Titolo XII  — Disposizioni finali tipizzate con numero manuale [N]
%
% RISULTATI ATTESI
%   - Il numero stampato in testa all'articolo corrisponde al valore
%     dell'argomento opzionale, indipendentemente dal contatore interno.
%   - Dopo un \articolo[N], il contatore automatico riprende dal valore
%     che aveva PRIMA dell'interruzione + 1 (non da N+1): effetto diretto
%     del \begingroup...\endgroup usato in \AN@emit.
%   - \artref{}, \cref{} e \Cref{} su un articolo numerato manualmente
%     restituiscono N, non il valore del contatore.
%   - Lo stesso comportamento vale per \comma[etichetta] e per le
%     disposizioni finali tipizzate con argomento opzionale.
% ===========================================================================

% ---------------------------------------------------------------------------
\titolo{Numerazione automatica — baseline di riferimento}
% ---------------------------------------------------------------------------
%
% Questi due articoli stabiliscono la baseline del contatore automatico
% subito prima dei test manuali.  I numeri effettivi dipendono da quanti
% articoli sono stati emessi in precedenza nel documento.

\articolo{Primo articolo automatico — baseline}[art:baseline1]
\comma{Articolo con numerazione automatica.  Il numero effettivo è
determinato dal contatore \texttt{subsubsection} al momento della
compilazione.  Verificare nell'indice degli articoli.}

\articolo{Secondo articolo automatico — baseline}[art:baseline2]
\comma{Deve avere il numero immediatamente successivo al precedente.
Sia $b$ tale numero: il primo articolo automatico dopo le interruzioni
manuali deve avere numero $b+1$.}

\notaredaz{%
  \textbf{TEST baseline.}\\
  Verificare che \artref{art:baseline1} e \artref{art:baseline2}
  abbiano numeri consecutivi nell'indice degli articoli e nel testo.
}

% ---------------------------------------------------------------------------
\titolo{Salto in avanti — \texttt{\textbackslash articolo[20]\{...\}}}
% ---------------------------------------------------------------------------

\articolo[20]{Articolo con numero manuale in avanti}[art:manuale20]
\comma{Il numero di questo articolo è stato forzato a~20 con
\texttt{\textbackslash articolo[20]\{...\}[art:manuale20]}.
Il contatore interno \textbf{non} viene modificato: l'articolo automatico
successivo riprenderà da $b+1$, non da~21.}[c:manuale20:1]

\notaredaz{%
  \textbf{TEST salto in avanti.}\\
  Verificare che:
  \begin{itemize}
    \item il titolo mostri «Articolo~20»;
    \item \artref{art:manuale20} $\to$ \textit{atteso:} art.~20;
    \item \cref{art:manuale20}   $\to$ \textit{atteso:} articolo~20;
    \item l'indice degli articoli riporti «Articolo~20».
  \end{itemize}
}

% ---------------------------------------------------------------------------
\titolo{Salto all'indietro — \texttt{\textbackslash articolo[5]\{...\}}}
% ---------------------------------------------------------------------------
%
% Il numero manuale è inferiore all'ultimo valore del contatore.
% La classe deve stampare 5 senza imporre vincoli di monotonia.

\articolo[5]{Articolo con numero manuale all'indietro}[art:manuale5]
\comma{Il numero di questo articolo è stato forzato a~5 anche se il
contatore interno ha già superato tale valore.
Questo caso verifica l'assenza di vincoli di monotonia nella
numerazione manuale.}

\notaredaz{%
  \textbf{TEST salto all'indietro.}\\
  Verificare che:
  \begin{itemize}
    \item il titolo mostri «Articolo~5» anche se il contatore è più alto;
    \item \artref{art:manuale5} $\to$ \textit{atteso:} art.~5;
    \item non vengano prodotti errori o warning di LaTeX.
  \end{itemize}
}

% ---------------------------------------------------------------------------
\titolo{Ripresa del contatore automatico dopo numerazione manuale}
% ---------------------------------------------------------------------------

\articolo{Articolo automatico dopo due numerazioni manuali}[art:dopo-manuale]
\comma{Questo articolo usa la numerazione automatica dopo
\artref{art:manuale20} (numero manuale~20) e \artref{art:manuale5}
(numero manuale~5).  Il suo numero deve essere $b+1$, dove $b$ è il
numero di \artref{art:baseline2}.  Se il numero fosse~21 o~6,
significa che \texttt{\textbackslash AN@emit} ha modificato il contatore
invece di isolarlo nel gruppo.}

\notaredaz{%
  \textbf{TEST ripresa contatore.}\\
  Sia $b$ = numero di \artref{art:baseline2}.
  Verificare che \artref{art:dopo-manuale} = $b+1$.
}

% ---------------------------------------------------------------------------
\titolo{Numero manuale con label e verifica riferimenti incrociati}
% ---------------------------------------------------------------------------

\articolo[42]{Articolo con numero manuale e label}[art:manuale42]
\comma{Articolo con numero forzato a~42 e label \texttt{art:manuale42}.
Tutti i comandi di riferimento devono restituire~42.}[c:manuale42:1]

\comma{Secondo comma dell'articolo~42.}[c:manuale42:2]

\notaredaz{%
  \textbf{TEST riferimenti su articolo con numero manuale.}\\
  \begin{itemize}
    \item \artref{art:manuale42}
          $\to$ \textit{atteso:} art.~42
    \item \articoloeref{art:manuale42}
          $\to$ \textit{atteso:} articolo~42
    \item \cref{art:manuale42}
          $\to$ \textit{atteso:} articolo~42
    \item \Cref{art:manuale42}
          $\to$ \textit{atteso:} Articolo~42
    \item \artcomma{art:manuale42}{c:manuale42:1}
          $\to$ \textit{atteso:} art.~42, comma~1
    \item \artcomma{art:manuale42}{c:manuale42:2}
          $\to$ \textit{atteso:} art.~42, comma~2
  \end{itemize}
}

% ---------------------------------------------------------------------------
\titolo{Comma con etichetta manuale e ripresa del contatore}
% ---------------------------------------------------------------------------

\articolo{Articolo con commi misti: automatici e manuali}[art:commi-misti]

\comma{Primo comma automatico — etichetta attesa: 1.}[c:commi-misti:1]

\comma{Secondo comma automatico — etichetta attesa: 2.}[c:commi-misti:2]

\comma[bis]{Comma con etichetta manuale «bis».
Il contatore \texttt{paragraph} \textbf{non} viene modificato
grazie al \texttt{\textbackslash begingroup}.}

\comma{Terzo comma automatico — etichetta attesa: 3, non 4.
Il «bis» non ha spostato il contatore.}

\comma[ter]{Altro comma con etichetta manuale «ter».}

\comma{Quarto comma automatico — etichetta attesa: 4.}

\notaredaz{%
  \textbf{TEST comma manuale e ripresa automatica.}\\
  Sequenza attesa per \artref{art:commi-misti}:
  1 $\to$ 2 $\to$ bis $\to$ 3 $\to$ ter $\to$ 4.\\
  Se la sequenza automatica fosse 1, 2, 3, 4, 5 (saltando le etichette
  manuali come se non esistessero) o si interrompesse, il
  \texttt{\textbackslash begingroup} non isola correttamente il contatore.
}

% ---------------------------------------------------------------------------
\titolo{Disposizioni finali tipizzate con numero manuale}
% ---------------------------------------------------------------------------
%
% \normatransitoria[N], \abrogazioni[N], \entrataInVigore[N]:
% il numero manuale deve comparire nel titolo e nell'indice degli articoli.

\normatransitoria[99]{Norma transitoria con numero manuale~99.
Il titolo deve mostrare «Articolo~99»
e l'indice deve riportare «Articolo~99 — Norma transitoria».}

\abrogazioni[100]{Abrogazioni con numero manuale~100.
Il titolo deve mostrare «Articolo~100»
e l'indice deve riportare «Articolo~100 — Abrogazioni».}

\entrataInVigore[101]{Entrata in vigore con numero manuale~101.
Il titolo deve mostrare «Articolo~101»
e l'indice deve riportare «Articolo~101 — Entrata in vigore».}

\notaredaz{%
  \textbf{TEST disposizioni finali tipizzate con [N].}\\
  Verificare che ciascuno dei tre articoli mostri il numero manuale
  nel titolo e nell'indice degli articoli, e che il contatore automatico
  non sia stato alterato (un eventuale \texttt{\textbackslash articolo}
  successivo senza argomento opzionale deve seguire la sequenza pre-interruzione).
}

% ===========================================================================
% INDICI FINALI
% ===========================================================================

\indiceallegati
\indicedefinizioni

% ===========================================================================
% SEZIONE H — Test accessibilità PDF/UA-1 (opzione accessible)
%
% Questa sezione è progettata per essere compilata con l'opzione
% "accessible" attiva (LuaLaTeX + \DocumentMetadata, vedere Sezione A).
%
% Il contenuto è visibile in entrambe le modalità (standard e accessible),
% ma il tagging PDF è presente solo quando l'opzione accessible è attiva.
%
% ISTRUZIONI:
%   1. Decommentare \DocumentMetadata e "accessible" in Sezione A.
%   2. Compilare con LuaLaTeX per tre passate.
%   3. Aprire il PDF in PAC 2024 o Adobe Acrobat Pro.
%   4. Verificare ciascun punto della checklist seguente.
% ===========================================================================

\titolo{Test di accessibilità PDF/UA-1}

% ---------------------------------------------------------------------------
% H.1 — Test: controllo prerequisiti e ClassWarning
% ---------------------------------------------------------------------------
%
% Il seguente blocco informativo spiega cosa verificare a livello
% di log di compilazione (non nel PDF).

\articolo{Verifica prerequisiti di compilazione}[art:acc-prerequisiti]

\comma{%
  La classe \texttt{attinormativi} v1.2 emette un \texttt{ClassWarning}
  nelle seguenti condizioni, verificabili nel log di compilazione
  (\texttt{.log}):
  \begin{lettere}
    \item opzione \texttt{accessible} attiva ma motore diverso da
          \hologo{LuaLaTeX}: il messaggio recita
          \emph{«L'opzione `accessible' richiede LuaLaTeX»};
    \item opzione \texttt{accessible} attiva ma
      \texttt{\textbackslash DocumentMetadata}
      assente: il messaggio recita
          \emph{«L'opzione `accessible' richiede \textbackslash DocumentMetadata
          prima di \textbackslash documentclass»}.
  \end{lettere}
  In entrambi i casi il tagging viene disattivato automaticamente
  e la compilazione prosegue senza errori fatali.%
}

\notaredaz{%
  \textbf{TEST H.1 — Prerequisiti e ClassWarning.}\\
  Verificare nel file \texttt{.log}:
  \begin{itemize}
    \item Con \texttt{accessible} attivo e LuaLaTeX + \texttt{\textbackslash DocumentMetadata}:
          nessun \texttt{ClassWarning} relativo all'accessibilità deve comparire.
    \item Ricompilare con \texttt{pdflatex} (mantenendo \texttt{accessible}):
          deve comparire il warning sul motore.
    \item Ricompilare con \texttt{lualatex} ma rimuovendo \texttt{\textbackslash DocumentMetadata}:
          deve comparire il warning su \texttt{\textbackslash DocumentMetadata}.
    \item In entrambi i casi il PDF viene prodotto senza errori e il
         booleano pubblico \texttt{\textbackslash ifANaccessible} vale \texttt{false}.
 \end{itemize}
}

% ---------------------------------------------------------------------------
% H.2 — Test: frontespizio taggato (Title + P)
% ---------------------------------------------------------------------------

\articolo{Frontespizio: strutture Title e P}[art:acc-frontespizio]

\comma{%
  In modalità \texttt{accessible}, il frontespizio generato da
  \texttt{\textbackslash frontespizioatto} (o dal suo alias
  \texttt{\textbackslash maketitle}) produce le seguenti strutture
  nell'albero dei tag PDF:
  \begin{lettere}
    \item il titolo ufficiale dell'atto (\texttt{\textbackslash titoloatto})
          è racchiuso in un tag \texttt{Title};
    \item gli altri campi (autorità, tipo atto, numero, data, oggetto)
          sono racchiusi ciascuno in un tag \texttt{P}.
  \end{lettere}%
}

\notaredaz{%
  \textbf{TEST H.2 — Frontespizio.}\\
  In PAC~2024 o nel pannello «Struttura» di Acrobat Pro, navigare
  alla prima pagina del documento e verificare:
  \begin{itemize}
    \item la presenza di un tag \texttt{Title} contenente il testo
          «Regolamento per la promozione delle energie rinnovabili
          e delle comunità energetiche»;
    \item la presenza di tag \texttt{P} distinti per ciascuno degli
          altri campi del frontespizio (autorità, tipo, numero, data, oggetto);
    \item l'assenza di contenuto non taggato (testo fuori struttura)
          sulla pagina del frontespizio.
  \end{itemize}
}

% ---------------------------------------------------------------------------
% H.3 — Test: gerarchia heading H1–H5
% ---------------------------------------------------------------------------

\articolo{Gerarchia strutturale: heading H1–H5}[art:acc-heading]

\comma{%
  La mappatura tra livelli normativi e ruoli PDF/UA è la seguente:
  \begin{numeri}
    \item \texttt{\textbackslash libro} $\to$ \texttt{H1};
    \item \texttt{\textbackslash titolo} $\to$ \texttt{H2};
    \item \texttt{\textbackslash capo} $\to$ \texttt{H3};
    \item \texttt{\textbackslash sezione} $\to$ \texttt{H4};
    \item \texttt{\textbackslash articolo} $\to$ \texttt{H5}.
  \end{numeri}
  Il comma (\texttt{\textbackslash comma}) non è un heading: il suo
  contenuto testuale è taggato come paragrafo (\texttt{P}), mentre
  il numero anteposto al testo è trattato come elemento decorativo.%
}

\notaredaz{%
  \textbf{TEST H.3 — Gerarchia heading.}\\
  In PAC~2024 (scheda «Logical Structure») o nel pannello
  «Ordine di lettura» di Acrobat Pro, verificare:
  \begin{itemize}
    \item i titoli di questo documento (\emph{Titolo~I},
          \emph{Titolo~II}, ecc.) risultano come \texttt{H2};
    \item i capi (\emph{Capo~I Benefici di carattere urbanistico},
          ecc.) risultano come \texttt{H3};
    \item gli articoli risultano come \texttt{H5};
    \item i commi risultano come \texttt{P}, non come \texttt{H6};
    \item la gerarchia è continua e non presenta salti di livello
          (es. \texttt{H2} seguito direttamente da \texttt{H5} senza
          \texttt{H3}/\texttt{H4} intermedi nello stesso ramo
          è corretto perché la classe non usa i livelli intermedi
          per i commi).
  \end{itemize}
}

% ---------------------------------------------------------------------------
% H.4 — Test: comma taggato come P
% ---------------------------------------------------------------------------

\articolo{Comma: tag P (paragrafo normativo)}[art:acc-comma]

\comma{Questo è il primo comma dell'articolo di test per il tag~P.
Il contenuto di questo comma deve apparire nell'albero dei tag come
paragrafo, non come heading.}[c:acc-comma:1]

\comma[bis]{Questo comma usa numerazione manuale «bis». Anche in questo
caso il contenuto deve essere taggato come \texttt{P}.}

\comma{Questo è il terzo comma, con numerazione automatica ripresa
correttamente dopo l'interruzione manuale.}[c:acc-comma:3]

\notaredaz{%
  \textbf{TEST H.4 — Comma come P.}\\
  Nell'albero dei tag, verificare per i commi di
  \artref{art:acc-comma}:
  \begin{itemize}
    \item ciascun comma è rappresentato da un tag \texttt{P};
    \item il numero anteposto (es. «1.», «bis.», «3.») è marcato come
          \texttt{Artifact} o non compare nell'albero dei tag come
          elemento testuale autonomo;
    \item la sequenza dei commi nell'ordine di lettura corrisponde
          all'ordine nel sorgente.
  \end{itemize}
}

% ---------------------------------------------------------------------------
% H.5 — Test: liste normative come OL
% ---------------------------------------------------------------------------

\articolo{Liste normative: tag OL/LI/LBody}[art:acc-liste]

\comma{Le liste di tipo \texttt{lettere} devono essere taggata come
lista ordinata. Esempio:
\begin{lettere}
  \item Prima voce della lista alfabetica;
  \item seconda voce della lista alfabetica;
  \item terza voce della lista alfabetica.
\end{lettere}%
}

\comma{Le liste di tipo \texttt{numeri} devono essere anch'esse
taggata come lista ordinata. Esempio:
\begin{numeri}
  \item Prima voce della lista numerica;
  \item seconda voce della lista numerica;
  \item terza voce della lista numerica.
\end{numeri}%
}

\notaredaz{%
  \textbf{TEST H.5 — Liste normative come OL.}\\
  Nell'albero dei tag, verificare per le liste di questo articolo:
  \begin{itemize}
    \item entrambe le liste risultano come \texttt{OL}
          (ordered list), non come \texttt{L} generica;
    \item ciascuna voce è una coppia \texttt{LI}/\texttt{LBody};
    \item l'etichetta alfabetica o numerica (es. «a)», «1)») è
          presente come \texttt{Lbl} all'interno di \texttt{LI};
    \item nessuna voce è fuori struttura.
  \end{itemize}
}

% ---------------------------------------------------------------------------
% H.6 — Test: ambiente definizioni come DL/DT/DD
% ---------------------------------------------------------------------------

\articolo{Ambiente definizioni: tag DL/DT/DD}[art:acc-definizioni]

\comma{Di seguito un sotto-insieme di definizioni ai soli fini
del test di accessibilità:
\begin{definizioni}
  \definizione{Termine alfa}{Prima definizione di test, usata per
    verificare la struttura \texttt{DT}/\texttt{DD}.}[def:acc-alfa]

  \definizione{Termine beta}{Seconda definizione di test. Il rimando
    al \defref{def:acc-alfa} deve essere navigabile.}[def:acc-beta]

  \definizione{Termine gamma}{Terza definizione di test, con testo
    più lungo per verificare che il tag \texttt{DD} contenga
    correttamente l'intera descrizione anche su più righe.}[def:acc-gamma]
\end{definizioni}
}

\notaredaz{%
  \textbf{TEST H.6 — Definizioni come DL/DT/DD.}\\
  Nell'albero dei tag, verificare:
  \begin{itemize}
    \item la lista è racchiusa in un tag \texttt{DL}
          (definition list);
    \item per ciascuna voce è presente una coppia:
          \texttt{DT} contenente il termine in grassetto
          e \texttt{DD} contenente il testo della definizione;
    \item la struttura si ripete per tutte e tre le voci
          senza interruzioni.
  \end{itemize}
}

% ---------------------------------------------------------------------------
% H.7 — Test: note redazionali come Artifact
% ---------------------------------------------------------------------------

\articolo{Note redazionali: tag Artifact}[art:acc-note]

\comma{In modalità \texttt{accessible}, le note redazionali (visibili
solo con l'opzione \texttt{bozza}) sono marcate come \texttt{Artifact}
nell'albero dei tag. I lettori di schermo non leggono gli artefatti,
poiché sono elementi decorativi o editoriali privi di valore normativo.%
\notaredazinline{questa nota inline deve risultare come Artifact nel PDF taggato}%
}

\notaredaz{%
  \textbf{TEST H.7 — Note redazionali come Artifact.}\\
  Nell'albero dei tag, verificare:
  \begin{itemize}
    \item questa nota a blocco (\texttt{\textbackslash notaredaz})
          non compare nell'albero dei tag come contenuto strutturato,
          oppure risulta esplicitamente marcata come \texttt{Artifact};
    \item la nota inline (\texttt{\textbackslash notaredazinline})
          del comma precedente è anch'essa marcata come \texttt{Artifact};
    \item il banner di bozza (\texttt{\textbackslash bozzainfo}),
          presente in apertura del documento, è marcato come
          \texttt{Artifact}.
  \end{itemize}
  In PAC~2024 il test «Artifacts are tagged correctly» deve risultare
  superato (verde) per questi elementi.
}

% ---------------------------------------------------------------------------
% H.8 — Test: preambolo e ambienti di servizio come Sect
% ---------------------------------------------------------------------------

\articolo{Ambienti di servizio: tag Sect}[art:acc-sect]

\comma{Il preambolo presente in apertura del documento (ambiente
\texttt{preambolo}) e l'eventuale relazione illustrativa (ambiente
\texttt{relazione}) sono racchiusi in un tag \texttt{Sect}
(sezione generica). Il tag \texttt{Sect} è il ruolo PDF/UA appropriato
per unità logiche dotate di un proprio titolo che non partecipano
alla gerarchia heading numerata dell'atto.%
}

\notaredaz{%
  \textbf{TEST H.8 — Preambolo e relazione come Sect.}\\
  Nell'albero dei tag, verificare:
  \begin{itemize}
    \item il blocco «Preambolo» in apertura del documento è racchiuso
          in un tag \texttt{Sect};
    \item all'interno della \texttt{Sect} è presente il titolo del
          capitolo non numerato («Preambolo»);
    \item il contenuto del preambolo (le clausole \emph{visto},
          \emph{ritenuto}, ecc.) è strutturato come paragrafi \texttt{P}
          all'interno della sezione.
  \end{itemize}
}

% ---------------------------------------------------------------------------
% H.9 — Test: allegati come Sect/Caption e indici come TOC
% ---------------------------------------------------------------------------

\articolo{Allegati e indici: tag Sect, Caption, TOC}[art:acc-allegati]

\comma{Gli allegati del documento sono strutturati come segue
nell'albero dei tag:
\begin{lettere}
  \item l'intero allegato è racchiuso in un tag \texttt{Sect};
  \item il titolo dell'allegato è marcato come \texttt{Caption}
        all'interno della \texttt{Sect};
  \item il contenuto dell'allegato è strutturato come paragrafi
        \texttt{P} e, se presenti, come articoli \texttt{H5} e
        commi \texttt{P}.
\end{lettere}%
}

\comma{Gli indici secondari (\texttt{\textbackslash indicearticoli},
\texttt{\textbackslash indiceallegati}, \texttt{\textbackslash indicedefinizioni})
generano ciascuno un blocco \texttt{TOC} (table of contents) nell'albero
dei tag. Il tag \texttt{TOC} è il ruolo PDF/UA dedicato agli indici
navigazionali.%
}

\notaredaz{%
  \textbf{TEST H.9 — Allegati e indici.}\\
  Per gli allegati (vedere Allegato~A e Allegato~B nel documento):
  \begin{itemize}
    \item navigare all'allegato nell'albero dei tag e verificare
          la struttura \texttt{Sect} $\to$ \texttt{Caption}
          (titolo) + contenuto;
    \item verificare che il titolo «Allegato~A --- Schema di statuto»
          risulti come \texttt{Caption}, non come \texttt{H2}.
  \end{itemize}
  Per gli indici (indicearticoli in apertura, indiceallegati e
  indicedefinizioni in chiusura):
  \begin{itemize}
    \item ciascun indice deve presentare un tag \texttt{TOC};
    \item le singole voci dell'indice devono risultare come
          \texttt{TOCI} (table of contents item) all'interno del
          \texttt{TOC}.
  \end{itemize}
}

% ---------------------------------------------------------------------------
% H.10 — Checklist PAC 2024 completa
% ---------------------------------------------------------------------------

\notaredaz{%
  \textbf{CHECKLIST ACCESSIBILITÀ v1.2 — Verifica PAC 2024}\\[0.5ex]
  Aprire il PDF in PAC~2024 e verificare che tutti i seguenti
  controlli risultino superati (icona verde):
  \begin{enumerate}
    \item \textbf{Document is tagged PDF} — il documento ha una
          struttura ad albero dei tag.
    \item \textbf{Document is marked as PDF/UA} — il dizionario
          \texttt{/MarkInfo} contiene \texttt{/Marked true} e
          il blocco \texttt{XMP} include il marcatore UA.
    \item \textbf{Document language is specified} — la lingua
          \texttt{it-IT} è dichiarata (da \texttt{\textbackslash DocumentMetadata}).
    \item \textbf{Title is specified} — il tag \texttt{Title}
          del frontespizio è presente.
    \item \textbf{Headings are used properly} — la gerarchia
          \texttt{H2}--\texttt{H5} è continua per ogni ramo.
    \item \textbf{Lists are structured properly} — le liste
          \texttt{lettere} e \texttt{numeri} risultano come
          \texttt{OL} con voci \texttt{LI}/\texttt{LBody}.
    \item \textbf{Artifacts are tagged correctly} — le note
          redazionali non compaiono come contenuto strutturato.
    \item \textbf{All content is contained in the structure tree} —
          nessun testo risulta fuori dall'albero dei tag.
  \end{enumerate}
  In caso di errori, annotare il numero del test PAC e il tag
  incriminato per il debug nella classe.
}

\end{document}
